\documentclass[../main.tex]{subfiles}
\begin{document}

\section{Проектирование агрегирующей базы данных}

Агрегирующая база данных предназначена для хранения снапшотов схем
всех источников. Она построена по нормализованной модели (третья
нормальная форма). Ключи, ограничения и индексы, в том числе
составные, хранятся реляционно: шапка объекта вынесена в одну
таблицу, состав колонок~--- в дочернюю, с сохранением порядка
следования. Это позволяет сравнивать описания структурно,
а не по тексту, и корректно обрабатывать составные ключи любой
длины.

\subsection{Модель данных}

Схема агрегирующей базы данных включает девять таблиц с общим
префиксом \field{sc\_} (schema catalog). Иерархия объектов
повторяет структуру самих СУБД: инстанс содержит схемы, схемы
содержат таблицы, таблицы содержат колонки, ограничения
и индексы. Ограничения и индексы, в свою очередь, имеют составные
колонки, вынесенные в отдельные таблицы. Одна дополнительная
таблица хранит историю федеративных запросов.

\subsubsection{Инстансы}

Таблица \field{sc\_instances} описывает источники, к которым
подключается сервис. Атрибуты: \field{id}~--- первичный ключ,
\field{label}~--- человекочитаемое имя,
\field{connection\_url}~--- JDBC-URL без учётных данных
(до 255 символов),
\field{username} и \field{password\_encrypted}~--- учётные данные,
пароль хранится в зашифрованном виде через функцию
\field{pgp\_sym\_encrypt},
\field{alias}~--- короткий уникальный идентификатор для
федеративных запросов,
\field{kind}~--- тип СУБД (POSTGRES, MARIADB, MYSQL, SQLITE,
MSSQL, OTHER),
\field{created\_at} и \field{updated\_at}~--- временные метки.
Поле \field{updated\_at} обновляется триггером
\field{sc\_instances\_set\_updated\_at} при каждом изменении строки.
На \field{label}, \field{connection\_url} и \field{alias} наложены
ограничения уникальности.

\subsubsection{Схемы и таблицы}

Таблица \field{sc\_schemas} хранит схемы (namespace) внутри
источника. Атрибуты: \field{id}, \field{instance\_id}~---
внешний ключ на инстанс, \field{schema\_name}, \field{owner},
\field{comment}. Пара \field{(instance\_id, schema\_name)} уникальна.

Таблица \field{sc\_tables} описывает объекты, у которых есть
колонки: обычные таблицы, партиционированные таблицы, представления,
материализованные представления и foreign tables. Атрибут
\field{kind} принимает одно из пяти значений. Пара
\field{(schema\_id, table\_name)} уникальна.

Таблица \field{sc\_columns} хранит колонки таблиц: имя,
порядковый номер (\field{ordinal\_position}), тип данных,
признак \field{is\_nullable}, значение по умолчанию
(\field{default\_value}), признак \field{is\_generated}
(GENERATED ALWAYS AS). Уникальность обеспечивается по двум парам:
\field{(table\_id, column\_name)} и \field{(table\_id, ordinal\_position)}.

\subsubsection{Ограничения}

Ограничения хранятся в двух таблицах. Таблица
\field{sc\_constraints} содержит \emph{шапку} ограничения:
\field{id}, \field{table\_id} (таблица, к которой относится
ограничение), \field{constraint\_name}, \field{type}
(PRIMARY\_KEY, FOREIGN\_KEY, UNIQUE, CHECK, EXCLUDE),
\field{expression} (тело для CHECK и EXCLUDE),
\field{ref\_table\_id} (целевая таблица для FOREIGN\_KEY),
\field{on\_delete} и \field{on\_update} (правила ссылочной
целостности для FOREIGN\_KEY: CASCADE, SET\_NULL, SET\_DEFAULT,
RESTRICT, NO\_ACTION).

Взаимоисключающие ограничения CHECK гарантируют корректность:
\field{ref\_table\_id} заполнен только для FOREIGN\_KEY,
\field{expression}~--- только для CHECK и EXCLUDE,
\field{on\_delete} и \field{on\_update}~--- только для FOREIGN\_KEY.

Таблица \field{sc\_constraint\_columns} содержит \emph{состав}
ограничения: одна строка на каждую колонку, входящую в ключ.
Атрибуты: \field{constraint\_id} (внешний ключ на
\field{sc\_constraints}), \field{column\_name}, \field{position}
(порядок колонки в ключе, начиная с единицы),
\field{ref\_column\_name} (для FOREIGN\_KEY~--- имя колонки
в целевой таблице, для остальных типов~--- NULL). Уникальность
по \field{(constraint\_id, column\_name)} и
\field{(constraint\_id, position)} исключает дублирование
и разрывы в порядке.

Такое разделение позволяет хранить составные ключи любой длины
и запрашивать их структуру без разбора текста. Например, для
составного первичного ключа таблицы \field{procuration\_authority}
вида \field{PRIMARY KEY (procuration\_id, authority\_id)}
в таблице \field{sc\_constraints} появится одна строка
с типом PRIMARY\_KEY, а в таблице \field{sc\_constraint\_columns}~---
две строки с позициями 1 и 2.

\subsubsection{Индексы}

Индексы устроены аналогично ограничениям. Таблица
\field{sc\_indexes} хранит шапку: \field{index\_name},
\field{is\_unique}, \field{is\_primary}, \field{kind} (BTREE,
HASH, GIN, GIST, SPGIST, BRIN) и \field{definition}~---
полный DDL, возвращаемый функцией \field{pg\_get\_indexdef()},
который необходим для выражений и частичных индексов.
Таблица \field{sc\_index\_columns} хранит состав индекса
с сохранением порядка: \field{column\_name}, \field{position}.

\subsubsection{История запросов}

Таблица \field{sc\_query\_history} хранит журнал федеративных
запросов: \field{sql\_text}, \field{status} (OK, ERROR, TIMEOUT,
CANCELLED), \field{row\_count}, \field{duration\_ms}, \field{error},
\field{created\_at}. Эта таблица изолирована: у неё нет внешних
ключей, потому что один федеративный запрос может обращаться
к нескольким источникам одновременно. Снимки содержимого строк
в агрегирующей базе не хранятся: сравнение строк выполняется
по запросу на клиенте, поэтому отдельная таблица для них не нужна.

\subsubsection{Диаграмма}

Диаграмма агрегирующей базы данных приведена
на рисунке~\ref{fig:schema-catalog}.

\begin{figure}[H]
  \centering
  \includegraphics[width=\linewidth]{\subfix{diagrams/schema-clean.png}}
  \caption{Схема агрегирующей базы данных (метаданные)}
  \label{fig:schema-catalog}
\end{figure}

\end{document}